\section{Антивирусные программы и методы защиты}
 
\subsection{Определение антивирусных программ}
 
Для обнаружения, удаления и защиты от компьютерных вирусов разработаны специальные программы, которые позволяют обнаруживать и уничтожать вирусы. Такие программы называются \textbf{антивирусными}. Современные антивирусные программы представляют собой многофункциональные продукты, сочетающие в себе как превентивные, профилактические средства, так и средства лечения вирусов и восстановления данных.
 
\subsection{Требования к антивирусным программам}
 
Количество и разнообразие вирусов велико, и чтобы их быстро и эффективно обнаружить, антивирусная программа должна отвечать некоторым параметрам.
 
\begin{equation} \label{eq:antivirus_quality}
    Q_{\text{AV}} = \frac{V_{\text{detected}}}{V_{\text{total}}} \times \alpha + \frac{U_{\text{updates}}}{U_{\text{required}}} \times \beta - \frac{F_{\text{false}}}{F_{\text{total}}} \times \gamma
\end{equation}
 
где:
\begin{itemize}
    \item \(Q_{\text{AV}}\) — качество антивирусной программы;
    \item \(V_{\text{detected}}\) — количество обнаруживаемых вирусов;
    \item \(V_{\text{total}}\) — общее количество известных вирусов;
    \item \(U_{\text{updates}}\) — частота обновлений;
    \item \(F_{\text{false}}\) — количество ложных срабатываний;
    \item \(\alpha, \beta, \gamma\) — весовые коэффициенты.
\end{itemize}
 
\subsection{Типы антивирусных программ}
 
Антивирусные программы делятся на:
 
\begin{enumerate}
    \item \textbf{Программы-детекторы} — обеспечивают поиск и обнаружение вирусов в оперативной памяти и на внешних носителях;
    \item \textbf{Программы-доктора (фаги)} — не только находят заражённые файлы, но и «лечат» их;
    \item \textbf{Программы-ревизоры} — запоминают исходное состояние программ и системных областей, затем сравнивают текущее состояние с исходным;
    \item \textbf{Программы-фильтры (сторожа)} — небольшие резидентные программы для обнаружения подозрительных действий;
    \item \textbf{Вакцины (иммунизаторы)} — резидентные программы, предотвращающие заражение файлов.
\end{enumerate}
 
\begin{table}[h]
    \centering
    \caption{Сравнительная характеристика типов антивирусных программ}
    \label{tab:antivirus_types}
    \begin{tabular}{|c|c|c|c|}
        \hline
        \textbf{Тип} & \textbf{Функция} & \textbf{Преимущества} & \textbf{Недостатки} \\
        \hline
        Детекторы & Поиск вирусов & Быстродействие & Не лечат файлы \\
        Доктора & Лечение файлов & Полное восстановление & Могут повредить файлы \\
        Ревизоры & Контроль изменений & Надёжность & Не удаляют вирусы \\
        Фильтры & Мониторинг действий & Раннее обнаружение & «Назойливость» \\
        Вакцины & Профилактика & Предотвращение заражения & Только от известных вирусов \\
        \hline
    \end{tabular}
\end{table}
 
\subsection{Краткий обзор антивирусных программ}
 
\subsubsection{Norton AntiVirus}
 
Один из наиболее известных и популярных антивирусов. Процент распознавания вирусов очень высокий (близок к 100\%). В программе используется механизм, позволяющий распознавать новые неизвестные вирусы. Антивирусные базы обновляются очень часто (иногда обновления появляются несколько раз в неделю).
 
\subsubsection{Dr. Web}
 
Популярный отечественный антивирус. Хорошо распознаёт вирусы, но в его базе их гораздо меньше, чем у других антивирусных программ.
 
\subsubsection{Antiviral Toolkit Pro (Лаборатория Касперского)}
 
Этот антивирус признан во всём мире как один из самых надёжных. Несмотря на простоту в использовании, он обладает всем необходимым арсеналом для борьбы с вирусами. Эвристический механизм, избыточное сканирование, сканирование архивов и упакованных файлов — это далеко не полный перечень его возможностей.
 
\subsection{Программная реализация: поиск сигнатур вирусов}
 
Ниже представлен фрагмент кода на Python, имитирующий простейший сигнатурный метод поиска вирусов:
 
\begin{listing}[h]
\caption{Сигнатурный метод обнаружения вирусов}
\label{code:signature_scanner}
 
\begin{minted}[frame=lines,linenos]{python}
import hashlib
import os
 
class SimpleAntivirus:
    """Упрощённая модель сигнатурного антивируса"""
    
    def __init__(self):
        self.signatures = {}  # словарь: MD5 сигнатуры вирусов
        self.virus_names = {}  # словарь: MD5 -> имя вируса
        
    def add_virus_signature(self, file_path, virus_name):
        """Добавление сигнатуры вируса в базу"""
        with open(file_path, 'rb') as f:
            content = f.read()
            md5_hash = hashlib.md5(content).hexdigest()
            self.signatures[md5_hash] = True
            self.virus_names[md5_hash] = virus_name
            print(f"Добавлена сигнатура вируса {virus_name}: {md5_hash}")
    
    def scan_file(self, file_path):
        """Сканирование файла на наличие вирусов"""
        try:
            with open(file_path, 'rb') as f:
                content = f.read()
                md5_hash = hashlib.md5(content).hexdigest()
                
                if md5_hash in self.signatures:
                    return f"ОБНАРУЖЕН ВИРУС: {self.virus_names[md5_hash]}"
                else:
                    return "Файл чист"
        except Exception as e:
            return f"Ошибка сканирования: {e}"
    
    def heuristic_scan(self, code):
        """Эвристическое сканирование (поиск подозрительных паттернов)"""
        suspicious_patterns = [
            b'DeleteFile', b'Format', b'WriteProcessMemory',
            b'CreateRemoteThread', b'RegSetValue'
        ]
        score = 0
        for pattern in suspicious_patterns:
            if pattern.lower() in code.lower():
                score += 20
                print(f"Обнаружен подозрительный паттерн: {pattern}")
        
        if score >= 60:
            return f"ВЫСОКИЙ РИСК: эвристический уровень {score}%"
        elif score >= 30:
            return f"СРЕДНИЙ РИСК: эвристический уровень {score}%"
        else:
            return f"НИЗКИЙ РИСК: эвристический уровень {score}%"
 
# Пример использования
av = SimpleAntivirus()
# av.add_virus_signature("virus_sample.exe", "TestVirus")
result = av.scan_file("clean_file.exe")
print(result)
\end{minted}
\end{listing}
 
\subsection{Пути проникновения вирусов}
 
Основными путями проникновения вирусов в компьютер являются:
 
\begin{enumerate}
    \item съёмные диски (гибкие и лазерные);
    \item компьютерные сети;
    \item электронная почта;
    \item загрузка файлов из интернета.
\end{enumerate}
 
Заражение жёсткого диска вирусами может произойти при загрузке программы с дискеты, содержащей вирус. Такое заражение может быть и случайным, например, если дискету не вынули из дисковода и перезагрузили компьютер.